\documentclass[%
fontsize=12pt, % Schriftgröße
twoside=off % kein einseitiges Layout
]{scrbook}
\usepackage{etoolbox}
\usepackage{scrlayer-scrpage} % Anpassbare Kopf- und Fußzeilen

\usepackage[utf8]{inputenc} % Textkodierung: UTF-8
\usepackage[T1]{fontenc} % Zeichensatzkodierung

\usepackage[english]{babel} % Deutsche Lokalisierung
\usepackage{graphicx} % Grafiken
\usepackage[absolute]{textpos} % Positionierung


\usepackage{calc} % Berechnungen
\usepackage{tabto} % Tabulatoren
\usepackage[hidelinks]{hyperref} % Hyperlinks
\usepackage[onehalfspacing]{setspace} % 1,5facher Zeilenabstand
\usepackage{calc} % Berechnungen
\usepackage{enumitem} % Mehr Kontrolle über itemize-, enumerate- und description-Umgebungen
\usepackage{relsize} % Schriftgröße in Abhängigkeit von aktueller anpassen
\usepackage{tabularx} % Flexiblere Tabellen
\usepackage[tablewithout, figurewithout]{caption} % Anpassen von Beschriftungen

\usepackage{bookmark} % Lesezeichen
\usepackage{amsmath}
\usepackage{mathtools}
\usepackage{amsfonts}
\usepackage{braket}

% Silbentrennung:
\usepackage{hyphenat}
\usepackage[
backend=biber,       % benutze biber statt bibtex
style=phys,          % Physikstil
citestyle=numeric,   % numerische Zitate [1], [2], ...
sorting=none,        % Reihenfolge wie im Dokument
natbib=true,         % optional für \citep, \citet
bibencoding=utf8,    % UTF-8 für .bib-Datei
hyperref=true,      % keine Links in Bibliographie
backref=false        % keine Rückverweise
]{biblatex}
\addbibresource{meta/literature.bib}

\newcommand{\hilbert}{\mathcal{H}}
\DeclareMathOperator{\Span}{span}
\begin{document}
\chapter{Static Yao-Kivelson Model}
\section{Kitaev Model and Yao-Kivelson Model}
Kitaev found a way to map the strongly-interacting spin model
\begin{equation}
H = - \sum_{\substack{\text{x-bonds} \\ \langle jk\rangle }} J^x_{jk} \sigma^x_j\sigma_x^k
- \sum_{\substack{\text{y-bonds} \\ \langle jk\rangle }} J^y_{jk} \sigma^y_j\sigma_y^k
- \sum_{\substack{\text{z-bonds} \\ \langle jk\rangle }} J^z_{jk} \sigma^z_j\sigma_z^k,
\label{eq:kitaevModel}
\end{equation}
now known as the Kitaev Model, to a free Majorana-Fermion hopping model, and therefore managed to solve the model exactly  \cite{kitaev_2006}. Classically, 'the ' Kitaev model is defined on a honeycomb lattice, where each coupling type $\alpha =x,y,z$ corresponds to one of the three bond orientations. However, his method directly generalizes to any tricoordinate lattice geometry \cite{reuther_2026}, provided that each site has exactly one bond of each type $\alpha = x,y,z$.

One alternative to the honeycomb lattice is the star/triangle-honeycomb lattice. Instead of elementary plaquettes of size 6, it features two different types of plaquettes of size 3 and 12, respectively. When defined on this geometry, model \ref{eq:kitaevModel} is known as the Yao-Kivelson Model \cite{yao_2007}. Due to the odd length of the triangular plaquettes, this model features a groundstate that spontaneously breaks Time-Reversal Symmetry, and therefore describes a Chiral Spin Liquid. 
\section{Kitaev's Spin-Majorana Mapping}
Kitaev introduced a representation of a spin-$\frac{1}{2}$-degree of freedom (called 'spin' in the following) in terms of 4 fermionic Majorana modes. Consider the 4-dimensional Fockspace $\mathcal{H}$ of two independent abstract fermions $a_1,a_2$ satisfying the canonical anti-commutation relations
\begin{equation}
	\Big\{a_k,a_l^\dagger\Big\}=\delta_{kl},\qquad \Big\{a_k,a_l\Big\}=0.
\end{equation}

They give rise to 4 real fermion modes
\begin{equation}
	c_{2k-1} \coloneq c_k + c_k^\dagger\quad\text{and}\quad c_{2k} \coloneq \frac{1}{i} \Big(a_k - a_k^\dagger\Big), \qquad\text{for}\qquad k=1,2,
	\label{eq:majoranasFromFermions}
\end{equation}
which are commonly called Majorana modes. They satisfy
\begin{equation}
	c_i^\dagger=c_i\quad\text{and}\quad\big\{c_i,c_j\big\}=2\delta_{ij} ,
	\label{eq:majoranaAntiCommutation}
\end{equation}
which qualifies them to be called "real fermions", in contrast to the usual complex ones.

The goal is now to use these to represent the spin algebra
\begin{equation}
	\sigma^\alpha \sigma^\beta
	= \delta^{\alpha\beta}
	+ i\epsilon^{\alpha\beta\gamma}\sigma^\gamma,
	\qquad \alpha,\beta,\gamma=x,y,z.
	\label{eq:spinAlgebra}
\end{equation}
For that, rename the 4 Majorana modes from $\{c_i\}_{i=1,\dots,4}$ to $\{b^x,b^y,b^z,d\}$ and define
\begin{equation}
	\tilde\sigma^\alpha \coloneq i b^\alpha d,\qquad\alpha=x,y,z.
	\label{eq:spinFromMajoranasDef}
\end{equation}
As $\left(\tilde\sigma^\alpha\right)^2=1$ follows directly from eq. \ref{eq:majoranaAntiCommutation}, eq. \ref{eq:spinAlgebra} reduces to the simpler, equivalent condition $\tilde\sigma^x\tilde\sigma^y\tilde\sigma^z\stackrel{!}{=}i$, which implies
\begin{equation}
	1 \stackrel{!}{=} b^x b^y b^z d \eqcolon D,
	\label{eq:dIdentityCondition}
\end{equation}
where we have defined the hermitean operator $D$. Due to $D^2 = 1$, $\mathcal{H}$ splits into the eigenspaces $\mathcal{H} = \mathcal{H}_1 \oplus \mathcal{H}_{-1}$, 
where the subscript denotes the respective eigenvalues of $D$. Consequently, $\mathcal{H}_1$ is the space where eq. $\ref{eq:dIdentityCondition}$ is fulfilled, and is therefore called the 'physical' subspace. The projector into this subspace is given by 
\begin{equation}
	P = \frac{1+D}{2}.
	\label{eq:projectorPhysical}
\end{equation}

D cannot be the identity on the whole Hilbert space $\mathcal{H}$, as it anti-commutes with each Majorana operator:
\begin{equation}
	\left\{D,c_i\right\} = 0\quad\text{for}\quad c_i = b^x,b^y,b^z,d.
	\label{eq:dAntiCommutesWithMajoranas}
\end{equation}


It follows that each Majorana operator $c_i$ maps 1-to-1 between $\hilbert_1$ and $\hilbert_{-1}$, which shows that $\dim(\hilbert_1) = \dim(\hilbert_{-1}) = 2$. Equation \ref{eq:dAntiCommutesWithMajoranas} further implies that $D$ commutes with operators consisting of an even number of Majorana operators. In particular,
\begin{equation}
	\left[D,\tilde\sigma^\alpha\right] = 0,\qquad\alpha=x,y,z,
	\label{eq:dCommutesWithSpins}
\end{equation}
which means $D$ and one (but not all) $\tilde\sigma^\alpha$, say $\tilde\sigma^z$, can be diagonalized simultaneously, such that we arrive at an orthonormal basis $\ket{D=\pm 1, \tilde\sigma^z=\pm 1}$ spanning $\mathcal{H}$:
\begin{equation}
	\mathcal{H}= \Span\{\underbrace{\ket{1,1},\ket{1,-1}}_\text{physical},\ket{-1,1},\ket{-1,-1}\}.
\end{equation}
This motivates a unitary transformation from the usual spinor space $\Span\{\ket{\uparrow},\ket{\downarrow}\}$ to the physical subspace $\mathcal{H}_1$ defined by
\begin{equation}
	U\ket{\uparrow} \coloneq \ket{1,1},\quad\text{and}\quad U\ket{\downarrow}\coloneq\ket{1,-1}.
\end{equation}
As the $\tilde\sigma^\alpha$ preserve both subspaces respectively due to eq. \ref{eq:dCommutesWithSpins}, their restriction to $\mathcal{H}_1$ finally defines a representation of the spin algebra sufficing
\begin{equation}
	\sigma^z = U^\dagger \tilde \sigma^z U,
\end{equation}
\begin{equation}
	[D,H] = [P,H]= 0
\end{equation}
for any physical Hamiltonian $H$ expressed in terms of the $\sigma^\alpha$ from eq. \ref{eq:spinFromMajoranasDef}. This is what allows for the following procedure to obtain eigenstates of a Hamiltonian:

\begin{enumerate}
	\item Write $H$ in terms of 4 Majorana modes per spin according to eq. \ref{eq:spinFromMajoranasDef}.
	\item Find the spectrum and eigenstates on the enlarged Hilbertspace.
	\item Apply $P$ to project into the physical subspace. The result will be an eigenstate of the physical Hamiltonian.
\end{enumerate}




\printbibliography
\end{document}